% Generated from locale/tr/content/second-order-logic/metatheory/undecidability-and-axiomatizability.tex; do not hand-edit.


\olfileid[tr]{sol}{met}{nax}

\olsection{İkinci Dereceden Mantık Aksiyomlaştırılamaz}


\begin{thm}
\ollabel{thm:sol-undecidable}
İkinci dereceden mantık karar verilemez.
\end{thm}

\begin{proof}
Birinci dereceden bir !!{sentence}, birinci dereceden mantıkta ancak ve ancak
ikinci dereceden mantıkta geçerlidir; birinci dereceden mantık karar
verilemez.
\end{proof}

\begin{thm}
\ollabel{cor:sol-not-axiomatizable}
İkinci dereceden mantık için sağlam ve tam bir !!{derivation} sistemi yoktur.
\end{thm}

\begin{proof}
$!A$, aritmetik dilinde !!a{sentence} olsun. $\Sat{N}{!A}$ ancak ve ancak
$\Th{PA^2}\Entails !A$. $!P$, $\Th{PA^2}$'nin dokuz aksiyomunun birleşimi
olsun. $\Th{PA^2}\Entails !A$ ancak ve ancak $\Entails !P\lif !A$, yani
her yapıda $\Sat{M}{!P\lif !A}$. Şimdi $\Obj{0}$ yerine $z$; $\prime$ yerine tek
konumlu fonksiyon değişkeni $u$; $+$ ve $\times$ yerine sırasıyla iki
konumlu fonksiyon değişkenleri $u'$ ile $u''$ ve $<$ yerine iki konumlu
bağıntı değişkeni $L$ koyup evrensel niceleyerek elde edilen
$\lforall[z][\lforall[u][\lforall[u'][\lforall[u''][\lforall[L][(!P' \lif !A')]]]]]$
!!{sentence} ifadesini ele alalım. Bu, saf ikinci dereceden mantığın geçerli bir
tümcesidir ancak ve ancak özgün tümce geçerliyse, ancak ve ancak
$\Th{PA^2}\Entails !A$ ise, ancak ve ancak $\Sat{N}{!A}$ ise. Dolayısıyla
ikinci dereceden mantık için sağlam ve tam bir ispat sistemi bulunsaydı,
$\Struct{N}$ içinde doğru olan !!{sentence}s kümesini numaralandıran
hesaplanabilir bir $f\colon\Nat\to\Sent[L_A]$ fonksiyonunu bu sistemden elde
edebilirdik. Bu fonksiyon $\Th{Q}$ içinde bir $!B_f(x,y)$ birinci dereceden
formülüyle temsil edilebilirdi. O zaman $\lexists[x][!B_f(x,y)]$
!!{formula} ifadesi, $\Struct{N}$ içinde doğru olan birinci dereceden
!!{sentence}s kümesini tanımlardı; bu da Tarski Teoremiyle çelişir.
\end{proof}



